\documentclass[11pt,a4paper]{article}

% ===================== PAQUETES =====================
\usepackage[utf8]{inputenc}
\usepackage[T1]{fontenc}
\usepackage[spanish]{babel}
\usepackage[margin=2.5cm]{geometry}
\usepackage{graphicx}
\usepackage{xcolor}
\usepackage{booktabs}
\usepackage{longtable}
\usepackage{listings}
\usepackage{tcolorbox}
\usepackage{hyperref}
\usepackage{enumitem}
\usepackage{fancyhdr}
\usepackage{titlesec}
\usepackage{amsmath}
\usepackage{amssymb}

% ===================== COLORES =====================
\definecolor{uteqazul}{RGB}{30,58,138}
\definecolor{codebg}{RGB}{245,245,245}
\definecolor{codegreen}{RGB}{0,128,0}
\definecolor{codegray}{RGB}{110,110,110}
\definecolor{codepurple}{RGB}{130,30,150}

% ===================== ESTILO DE CÓDIGO =====================
\lstdefinestyle{javastyle}{
    backgroundcolor=\color{codebg},
    commentstyle=\color{codegreen},
    keywordstyle=\color{uteqazul}\bfseries,
    stringstyle=\color{codepurple},
    numberstyle=\tiny\color{codegray},
    basicstyle=\ttfamily\footnotesize,
    breakatwhitespace=false,
    breaklines=true,
    captionpos=b,
    keepspaces=true,
    numbers=left,
    numbersep=5pt,
    showspaces=false,
    showstringspaces=false,
    showtabs=false,
    tabsize=2,
    frame=single,
    rulecolor=\color{black!20},
    language=Java
}
\lstset{style=javastyle}

% ===================== ENCABEZADOS =====================
\pagestyle{fancy}
\fancyhf{}
\fancyhead[L]{\small Aplicaciones Distribuidas (ISR-701) — EV-PR-U2}
\fancyhead[R]{\small MercadoQuevedo}
\fancyfoot[C]{\thepage}
\renewcommand{\headrulewidth}{0.4pt}

\titleformat{\section}{\normalfont\Large\bfseries\color{uteqazul}}{\thesection}{1em}{}
\titleformat{\subsection}{\normalfont\large\bfseries\color{uteqazul}}{\thesubsection}{1em}{}

\hypersetup{
    colorlinks=true,
    linkcolor=uteqazul,
    urlcolor=uteqazul,
    citecolor=uteqazul,
    pdftitle={MercadoQuevedo - Examen Practico Unidad 2},
    pdfauthor={Harold Vinueza}
}

\newcommand{\ok}{\textcolor{green!50!black}{\checkmark}}

% ===================== DOCUMENTO =====================
\begin{document}

% ---------- PORTADA ----------
\begin{titlepage}
    \centering
    \vspace*{1cm}
    {\large UNIVERSIDAD TÉCNICA ESTATAL DE QUEVEDO\par}
    {\large Facultad de Ciencias de la Computación y Diseño Digital\par}
    {\large Carrera de Ingeniería de Software\par}
    \vspace{1.5cm}

    \rule{\linewidth}{0.5mm}\\[0.4cm]
    {\huge\bfseries Examen Práctico — Unidad 2\par}
    \vspace{0.3cm}
    {\Large Arquitecturas Distribuidas: Microservicios, APIs RESTful\\ y Contenerización\par}
    \vspace{0.2cm}
    {\large Caso: MercadoQuevedo — Marketplace distribuido\par}
    \rule{\linewidth}{0.5mm}\\[1.5cm]

    \begin{tabbing}
        \hspace{4cm} \= \kill
        \textbf{Asignatura:} \> Aplicaciones Distribuidas (ISR-701) \\[0.2cm]
        \textbf{Unidad:} \> 2 — Arquitecturas Distribuidas \\[0.2cm]
        \textbf{Código del instrumento:} \> EV-PR-U2 \\[0.2cm]
        \textbf{Modalidad:} \> Individual — Laboratorio con defensa técnica oral \\[0.2cm]
        \textbf{Estudiante:} \> Harold Vinueza \\[0.2cm]
        \textbf{Docente:} \> Ing. Guerrero Ulloa Gleiston Cicerón, Mg.Sc. \\[0.2cm]
        \textbf{Período académico:} \> 2026--2027 \\[0.2cm]
        \textbf{Repositorio:} \> \url{https://github.com/Harold-Vinueza/EV-SUM-02-Examen-Parcial-I-Unidades-1-y-2} \\[0.2cm]
        \textbf{Stack:} \> Java 21 · Spring Boot 3.3 · PostgreSQL 16 · Docker Compose
    \end{tabbing}

    \vfill
    {\large Quevedo, Los Ríos, Ecuador — Julio 2026\par}
\end{titlepage}

\tableofcontents
\newpage

% ======================================================
\section{Introducción}
% ======================================================

El presente informe documenta el diseño e implementación de \textbf{MercadoQuevedo},
un prototipo de arquitectura de microservicios desarrollado para el Examen
Práctico de la Unidad 2 de la asignatura Aplicaciones Distribuidas (ISR-701).
El sistema resuelve el caso planteado por el docente: una plataforma de
marketplace para comerciantes locales del cantón Quevedo, migrando desde un
monolito conceptual hacia una arquitectura de microservicios que evidencia
separación de dominios, diseño RESTful disciplinado y despliegue reproducible
mediante contenedores.

El prototipo implementa dos microservicios independientes —\texttt{svc-catalogo}
y \texttt{svc-pedidos}— siguiendo el patrón \emph{Database per Service}
\cite{richardson2018}, comunicados de forma síncrona vía REST con manejo
explícito de fallos, contenerizados con Docker mediante \emph{multi-stage
builds}, y orquestados con Docker Compose detrás de un \emph{reverse proxy}
Nginx.

Este documento se organiza siguiendo las cinco partes de la rúbrica de
evaluación (C1--C7): arquitectura y diseño en capas, diseño de la API REST,
comunicación entre servicios, contenerización, orquestación, calidad del
repositorio, y el componente de análisis y defensa técnica.

% ======================================================
\section{Parte A — Arquitectura de Microservicios y Diseño en Capas}
% ======================================================

\subsection{Descomposición del dominio}

Siguiendo el enunciado del examen, el dominio de MercadoQuevedo se descompuso
en dos capacidades de negocio independientes, cada una con su propio
microservicio y su propia base de datos (\emph{Database per Service}
\cite{richardson2018}):

\begin{itemize}[leftmargin=2em]
    \item \textbf{\texttt{svc-catalogo}}: altas, bajas, modificaciones y
    consultas de productos (SKU, nombre, precio, stock).
    \item \textbf{\texttt{svc-pedidos}}: creación de pedidos que referencian
    uno o más productos, cálculo de subtotales y consulta del estado del pedido.
\end{itemize}

Ninguno de los dos servicios comparte esquema ni tabla de base de datos con el
otro; \texttt{svc-pedidos} solo conserva una \emph{referencia externa}
(\texttt{productoId}) al recurso remoto, nunca una copia de su estado interno.

\subsection{Diagrama de arquitectura}

La Figura~\ref{fig:arquitectura} muestra el diagrama de despliegue (UML / C4
Nivel 2) del sistema. El diagrama fuente editable se encuentra en
\texttt{docs/architecture.drawio} del repositorio (formato draw.io/diagrams.net).

\begin{figure}[h!]
    \centering
    \includegraphics[width=0.95\textwidth]{figures/architecture.png}
    \caption{Diagrama de despliegue de MercadoQuevedo. Cliente $\to$ Nginx
    (única puerta de entrada, puerto 8080) $\to$ \texttt{svc-catalogo} /
    \texttt{svc-pedidos}, cada uno con su base de datos PostgreSQL propia.
    La flecha punteada representa la comunicación síncrona inter-servicio
    con timeout explícito.}
    \label{fig:arquitectura}
\end{figure}

\textbf{Nota de compilación (Overleaf):} exporta \texttt{docs/architecture.drawio}
como PNG (draw.io $\to$ File $\to$ Export as $\to$ PNG, resolución 300 DPI) y
súbelo a \texttt{figures/architecture.png} en tu proyecto de Overleaf antes de
compilar.

\subsection{Diseño en capas}

Cada microservicio se organiza en cuatro capas con responsabilidades
estrictamente separadas, siguiendo el patrón arquitectónico visto en la
Unidad 2 (arquitectura en capas dentro de cada microservicio):

\begin{table}[h!]
    \centering
    \begin{tabular}{@{}ll@{}}
        \toprule
        \textbf{Capa} & \textbf{Responsabilidad} \\
        \midrule
        \texttt{controller} & Expone los endpoints REST, traduce HTTP $\leftrightarrow$ DTOs \\
        \texttt{service} & Reglas de negocio (SKU único, stock, existencia) \\
        \texttt{repository} & Acceso a datos vía Spring Data JPA \\
        \texttt{model} & Entidades JPA (\texttt{Producto}, \texttt{Pedido}, \texttt{PedidoItem}) \\
        \texttt{dto} & Contratos de entrada/salida (\emph{records} inmutables) \\
        \texttt{exception} & Excepciones de dominio + \texttt{@RestControllerAdvice} \\
        \bottomrule
    \end{tabular}
    \caption{Capas dentro de cada microservicio}
\end{table}

El \texttt{controller} nunca accede directamente al \texttt{repository}, y la
entidad JPA (\texttt{Producto}) nunca se serializa directamente al cliente —
siempre se traduce a un DTO (\texttt{ProductoResponse}), evitando acoplar el
contrato público de la API al esquema interno de persistencia.

% ======================================================
\section{Parte B — Diseño de la API REST}
% ======================================================

\subsection{Nivel de madurez de Richardson}

La API se ubica en el \textbf{Nivel 2} del modelo de madurez de Richardson
\cite{richardson2018}: los recursos se identifican por URI, se usan verbos
HTTP de forma disciplinada, y los códigos de estado son semánticamente
correctos. No se alcanza el Nivel 3 porque no se implementa hipermedia
(HATEOAS) más allá del header \texttt{Location}.

\subsection{Convenciones de URI}

Todas las URIs siguen las convenciones REST exigidas: sustantivos en plural,
versionadas (\texttt{/api/v1/...}), sin verbos en la ruta, con los
identificadores en el \emph{path} (nunca en \emph{query string}).

\begin{table}[h!]
    \centering
    \small
    \begin{tabular}{@{}lllll@{}}
        \toprule
        \textbf{Recurso} & \textbf{Método} & \textbf{URI} & \textbf{Éxito} & \textbf{Errores} \\
        \midrule
        Productos & POST   & \texttt{/api/v1/productos}      & 201 + \texttt{Location} & 400, 409, 422 \\
        Productos & GET    & \texttt{/api/v1/productos}      & 200 & --- \\
        Productos & GET    & \texttt{/api/v1/productos/\{id\}} & 200 & 404 \\
        Productos & PUT    & \texttt{/api/v1/productos/\{id\}} & 200 (idemp.) & 404, 409, 422 \\
        Productos & DELETE & \texttt{/api/v1/productos/\{id\}} & 204 (idemp.) & 404 \\
        Pedidos   & POST   & \texttt{/api/v1/pedidos}        & 201 + \texttt{Location} & 400, 404, 409, 422, 503 \\
        Pedidos   & GET    & \texttt{/api/v1/pedidos}        & 200 & --- \\
        Pedidos   & GET    & \texttt{/api/v1/pedidos/\{id\}} & 200 & 404 \\
        \bottomrule
    \end{tabular}
    \caption{Endpoints expuestos por \texttt{svc-catalogo} y \texttt{svc-pedidos}}
    \label{tab:endpoints}
\end{table}

\subsection{Idempotencia}

\texttt{PUT} y \texttt{DELETE} son idempotentes por diseño: invocar
\texttt{PUT /api/v1/productos/1} múltiples veces con el mismo cuerpo produce
siempre el mismo estado final; \texttt{DELETE} sobre un recurso ya eliminado
devuelve \texttt{404} en la segunda invocación, pero el \emph{estado del
sistema} no cambia entre la primera y las siguientes llamadas —justamente lo
que exige la definición de idempotencia. \texttt{POST}, en cambio, no es
idempotente: cada invocación crea un nuevo recurso.

\subsection{Manejo de errores: distinción 400 vs.\ 422}

Un aspecto cuidado del diseño es la distinción exigida por la rúbrica entre
\texttt{400 Bad Request} (JSON sintácticamente inválido) y
\texttt{422 Unprocessable Entity} (JSON válido que viola una regla de
dominio). En Spring, esto se logra capturando dos excepciones distintas en el
\texttt{GlobalExceptionHandler}:

\begin{lstlisting}[caption={Distinción entre JSON malformado (400) y validación de dominio (422)}, label={lst:exceptionhandler}]
@ExceptionHandler(MethodArgumentNotValidException.class)
public ResponseEntity<ErrorResponse> manejarValidacion(
        MethodArgumentNotValidException ex) {
    // @Valid fallo una regla de negocio -> JSON SI era valido
    String detalle = ex.getBindingResult().getFieldErrors().stream()
            .map(err -> err.getField() + ": " + err.getDefaultMessage())
            .collect(Collectors.joining("; "));
    return ResponseEntity.status(HttpStatus.UNPROCESSABLE_ENTITY)
            .body(new ErrorResponse(detalle, "VALIDACION_FALLIDA"));
}

@ExceptionHandler(HttpMessageNotReadableException.class)
public ResponseEntity<ErrorResponse> manejarJsonMalformado(
        HttpMessageNotReadableException ex) {
    // Jackson ni siquiera pudo parsear el JSON -> ocurre ANTES de @Valid
    return ResponseEntity.status(HttpStatus.BAD_REQUEST)
            .body(new ErrorResponse(
                "El cuerpo de la peticion no es JSON valido",
                "JSON_MALFORMADO"));
}
\end{lstlisting}

\subsection{Documentación OpenAPI}

Cada microservicio publica su especificación OpenAPI 3.0 mediante
\texttt{springdoc-openapi-starter-webmvc-ui}, accesible en
\texttt{/api/v1/openapi.json} y visualizable en Swagger UI
(\texttt{/api/v1/docs}). Los artefactos exportados se encuentran en
\texttt{docs/api/catalogo.json} y \texttt{docs/api/pedidos.json} del
repositorio.

% ======================================================
\section{Parte C — Comunicación entre Servicios}
% ======================================================

\subsection{Comunicación síncrona con timeout explícito}

Al crear un pedido, \texttt{svc-pedidos} invoca por HTTP a
\texttt{svc-catalogo} para verificar existencia y disponibilidad del producto.
La comunicación se implementa con \texttt{RestClient} (Spring Framework 6.1+)
configurado con un \texttt{connectTimeout}/\texttt{readTimeout} explícito de
2 segundos:

\begin{lstlisting}[caption={Cliente HTTP con timeout explícito hacia \texttt{svc-catalogo}}, label={lst:catalogoclient}]
SimpleClientHttpRequestFactory factory =
        new SimpleClientHttpRequestFactory();
int timeoutMillis = (int) (timeoutSeconds * 1000);
factory.setConnectTimeout(timeoutMillis);
factory.setReadTimeout(timeoutMillis);

this.restClient = RestClient.builder()
        .baseUrl(baseUrl)
        .requestFactory(factory)
        .build();
\end{lstlisting}

\subsection{Manejo de fallos}

Tres escenarios de fallo se manejan explícitamente, cada uno mapeado a un
código HTTP semánticamente correcto (Tabla~\ref{tab:fallos}):

\begin{table}[h!]
    \centering
    \small
    \begin{tabular}{@{}p{4.2cm}p{2cm}p{6.5cm}@{}}
        \toprule
        \textbf{Escenario} & \textbf{Código} & \textbf{Mecanismo} \\
        \midrule
        Producto no existe en el catálogo & 404 & \texttt{HttpClientErrorException.NotFound} $\to$ \texttt{ProductoInexistenteException} \\
        Stock insuficiente & 409 & Regla de negocio en \texttt{PedidoService}, \texttt{StockInsuficienteException} \\
        \texttt{svc-catalogo} caído / colgado & 503 & \texttt{ResourceAccessException} (timeout o conexión rechazada) $\to$ \texttt{CatalogoNoDisponibleException} \\
        \bottomrule
    \end{tabular}
    \caption{Manejo de fallos en la comunicación inter-servicio}
    \label{tab:fallos}
\end{table}

\begin{lstlisting}[caption={Captura y traducción de fallos de red a excepciones de dominio}, label={lst:catchfallos}]
try {
    return restClient.get()
            .uri("/api/v1/productos/{id}", productoId)
            .retrieve()
            .body(ProductoDTO.class);
} catch (HttpClientErrorException.NotFound e) {
    throw new ProductoInexistenteException(productoId);
} catch (HttpServerErrorException e) {
    throw new CatalogoNoDisponibleException(
            "svc-catalogo respondio con error del servidor: "
            + e.getStatusCode());
} catch (ResourceAccessException e) {
    // Cubre timeout y conexion rechazada
    throw new CatalogoNoDisponibleException(
            "svc-catalogo no respondio en " + timeoutSeconds
            + "s: " + e.getMessage());
}
\end{lstlisting}

Nunca se deja que el error de red se propague como un \texttt{500} genérico,
ni que la petición se cuelgue indefinidamente esperando al catálogo —esto es
central para la pregunta 6 del componente de análisis (Sección~\ref{sec:analisis}).

% ======================================================
\section{Parte D — Contenerización con Docker}
% ======================================================

\subsection{\emph{Multi-stage build}}

Cada microservicio cuenta con un \texttt{Dockerfile} de dos etapas: una etapa
de compilación (con el JDK completo y Maven) y una etapa de \emph{runtime}
(solo el JRE), siguiendo las buenas prácticas documentadas por
Karabey~Aksakalli~et~al.~\cite{aksakalli2021}.

\begin{lstlisting}[language={}, basicstyle=\ttfamily\footnotesize, caption={Dockerfile multi-stage de \texttt{svc-catalogo}}, label={lst:dockerfile}]
# ---------- Etapa 1: build ----------
FROM maven:3.9-eclipse-temurin-21 AS build
WORKDIR /build
COPY pom.xml .
RUN mvn -B dependency:go-offline
COPY src ./src
RUN mvn -B clean package -DskipTests

# ---------- Etapa 2: runtime ----------
FROM eclipse-temurin:21-jre-alpine AS runtime
RUN apk add --no-cache curl \
    && addgroup -S app && adduser -S app -G app
WORKDIR /app
COPY --from=build /build/target/svc-catalogo.jar app.jar
USER app
EXPOSE 8000
HEALTHCHECK --interval=10s --timeout=3s --start-period=20s --retries=5 \
    CMD curl -f http://localhost:8000/health || exit 1
ENTRYPOINT ["java", "-jar", "app.jar"]
\end{lstlisting}

\subsection{Justificación de la imagen base}

\begin{itemize}[leftmargin=2em]
    \item \textbf{Etapa build} (\texttt{maven:3.9-eclipse-temurin-21}): incluye
    JDK completo + Maven, necesarios solo para compilar. Se descarta por completo
    de la imagen final.
    \item \textbf{Etapa runtime} (\texttt{eclipse-temurin:21-jre-alpine}): solo
    el JRE sobre Alpine Linux —sin compilador, sin Maven, sin código fuente—,
    resultando en una imagen de aproximadamente 180--220~MB frente a los
    500+~MB de la imagen de construcción, con una superficie de ataque mínima.
    \item Ambas imágenes ejecutan como \textbf{usuario no-root} (\texttt{app}),
    creado explícitamente en el \texttt{Dockerfile}, y exponen únicamente el
    puerto \texttt{8000}.
\end{itemize}

% ======================================================
\section{Parte E — Orquestación con Docker Compose}
% ======================================================

El archivo \texttt{docker-compose.yml} orquesta cinco servicios: las dos
bases de datos PostgreSQL, los dos microservicios y Nginx como puerta de
entrada única.

\begin{lstlisting}[language={}, basicstyle=\ttfamily\footnotesize, caption={Extracto de \texttt{docker-compose.yml}: healthchecks encadenados}, label={lst:compose}]
svc-pedidos:
  build: ./svc-pedidos
  depends_on:
    db-pedidos:
      condition: service_healthy
    svc-catalogo:
      condition: service_healthy
  healthcheck:
    test: ["CMD", "curl", "-f", "http://localhost:8000/health"]
    interval: 10s
    timeout: 3s
    retries: 8
    start_period: 30s
\end{lstlisting}

\textbf{Características clave:}
\begin{itemize}[leftmargin=2em]
    \item Red \texttt{mercadoquevedo-net} (\emph{bridge}) dedicada: ningún
    servicio interno tiene puertos publicados al host, salvo Nginx
    (\texttt{8080:80}).
    \item \texttt{depends\_on} con \texttt{condition: service\_healthy} en
    cascada: las bases de datos deben estar \emph{healthy} antes de que
    arranquen los microservicios, y éstos deben estarlo antes de que Nginx los
    reciba como \emph{upstream}.
    \item Volúmenes nombrados (\texttt{catalogo\_data}, \texttt{pedidos\_data})
    para persistencia entre reinicios.
    \item Nginx enruta \texttt{/api/v1/productos/*} hacia \texttt{svc-catalogo}
    y \texttt{/api/v1/pedidos/*} hacia \texttt{svc-pedidos}; es el único
    servicio con healthcheck basado en \texttt{curl} sobre una imagen propia
    (\texttt{nginx/Dockerfile}) que agrega dicha dependencia, ya que la imagen
    oficial \texttt{nginx:alpine} reciente no la incluye por defecto.
\end{itemize}

% ======================================================
\section{Parte F — Calidad del Código y del Repositorio}
% ======================================================

El repositorio se organiza siguiendo la convención estándar de un proyecto
Maven multi-módulo informal (dos proyectos independientes en carpetas
hermanas), con:

\begin{itemize}[leftmargin=2em]
    \item Historial de commits descriptivo por funcionalidad.
    \item \texttt{README.md} con instrucciones completas de compilación y
    despliegue (\texttt{docker compose up --build}, un solo comando).
    \item Especificación OpenAPI publicada en \texttt{docs/api/}.
    \item \texttt{.gitignore} y \texttt{.dockerignore} excluyendo artefactos
    de build (\texttt{target/}, \texttt{*.class}) del control de versiones y
    del contexto de construcción de Docker.
    \item \texttt{ANALISIS.md} con el componente de análisis y defensa técnica
    (reproducido y ampliado en la Sección~\ref{sec:analisis}).
\end{itemize}

% ======================================================
\section{Componente de Análisis y Defensa Técnica}
\label{sec:analisis}
% ======================================================

\subsection{¿Cuál restricción REST de Fielding identifica con mayor claridad?}

La más clara es la \textbf{interfaz uniforme} (\emph{uniform interface})
\cite{fielding2000}, en particular sus sub-restricciones de \emph{identificación
de recursos} y \emph{manipulación de recursos a través de representaciones}.
En \texttt{ProductoController}, el identificador \texttt{\{id\}} localiza el
recurso solo por URI (sin verbos en la ruta), y el cliente lo manipula
únicamente a través de representaciones (\emph{records} \texttt{ProductoRequest}/
\texttt{ProductoResponse}), nunca accediendo a la entidad JPA directamente. El
\texttt{ResponseEntity.created(URI...)} es además un mensaje autodescriptivo
que indica al cliente dónde vive el nuevo recurso. También se cumple
\textbf{stateless}: ningún controlador guarda sesión ni contexto entre
peticiones.

\subsection{¿En qué nivel del modelo de madurez de Richardson se ubica la API?}

\textbf{Nivel 2}: recursos identificados por URI, verbos HTTP disciplinados,
códigos de estado semánticamente correctos, pero sin hipermedia (HATEOAS) más
allá del header \texttt{Location} —lo que impide alcanzar el Nivel 3. Ejemplo
concreto: \texttt{POST /api/v1/productos} con SKU repetido devuelve
\texttt{409 Conflict} con cuerpo \texttt{\{"codigo":"SKU\_DUPLICADO"\}}, en
lugar de un genérico \texttt{200 OK} que sería típico de un Nivel 0/1.

\subsection{Acoplamiento residual entre \texttt{svc-catalogo} y \texttt{svc-pedidos} (Newman)}

Pese al desacoplamiento de datos (\emph{Database per Service}), Newman
\cite{newman2021} señala que \emph{coupling} y \emph{cohesion} deben
evaluarse más allá del esquema de datos. Se identifican dos acoplamientos
residuales:

\begin{enumerate}[leftmargin=2em]
    \item \textbf{Acoplamiento temporal}: \texttt{PedidoService.crear()}
    llama a \texttt{CatalogoClient} de forma síncrona y bloqueante; ambos
    servicios deben estar vivos simultáneamente.
    \item \textbf{Acoplamiento de contrato}: \texttt{svc-pedidos} deserializa
    directamente la respuesta de \texttt{svc-catalogo} en \texttt{ProductoDTO};
    un cambio de esquema no versionado rompería la deserialización en tiempo
    de ejecución, sin aviso en compilación.
\end{enumerate}

\textbf{Mitigación en producción}: migrar la verificación de stock a un
modelo orientado a eventos (EDA) con Kafka/RabbitMQ (patrón CQRS, réplica de
solo lectura en \texttt{svc-pedidos}); introducir pruebas de contrato
(Spring Cloud Contract/Pact) sobre el OpenAPI publicado; y versionar la API
(\texttt{/api/v1/} $\to$ \texttt{/api/v2/}) ante cambios incompatibles.

\subsection{Justificación de microservicios pese al hallazgo de Blinowski et al.}

Blinowski et al.~\cite{blinowski2022} muestran que, sobre una sola máquina,
un monolito puede superar en rendimiento a su equivalente en microservicios,
por el \emph{overhead} de red, serialización y arranque de múltiples JVMs.
Ese hallazgo es correcto pero mide \emph{rendimiento crudo}, no las razones
reales de migración en MercadoQuevedo: escalado independiente ante demanda
diferenciada entre catálogo y pedidos, aislamiento de fallos (demostrado en
el prototipo: si \texttt{svc-catalogo} cae, \texttt{svc-pedidos} responde
\texttt{503} controlado en vez de caer con él), y autonomía organizacional a
futuro cuando existan varios equipos. Como advierte el propio Newman, la
complejidad operativa de los microservicios ``solo se paga cuando sus
beneficios se necesitan''; aquí se paga deliberadamente como prueba de
concepto, no por necesidad de throughput actual.

\subsection{REST vs.\ gRPC entre \texttt{svc-pedidos} y \texttt{svc-catalogo}}

\textbf{Se ganaría}: menor latencia y tamaño de payload (Protocol Buffers
binario sobre HTTP/2 vs.\ JSON sobre HTTP/1.1), tipado fuerte del contrato en
tiempo de compilación vía \texttt{.proto} (eliminando el riesgo de
deserialización silenciosa descrito arriba), y \emph{streaming} nativo.
\textbf{Se perdería}: accesibilidad para consumidores externos —el examen
exige que \texttt{svc-catalogo} sea consumible con \texttt{curl}/Postman y
documentado en OpenAPI/Swagger, algo que gRPC no ofrece nativamente sin
herramientas especializadas (\texttt{grpcurl})— y la simplicidad de depuración
durante la demostración en vivo. \textbf{Conclusión}: mantener REST público
en \texttt{svc-catalogo} y evaluar gRPC únicamente para la llamada interna si
el volumen de tráfico lo justificara.

\subsection{Si \texttt{svc-catalogo} se detiene: código HTTP y Circuit Breaker}

\textbf{Código recibido}: \texttt{503 Service Unavailable}
(\texttt{CATALOGO\_NO\_DISPONIBLE}), porque el \texttt{ConnectException}
(conexión rechazada) o el \texttt{SocketTimeoutException} (tras 2s) se
envuelven en \texttt{ResourceAccessException}, capturada explícitamente y
traducida a \texttt{CatalogoNoDisponibleException}, mapeada a \texttt{503} por
el \texttt{GlobalExceptionHandler}.

\textbf{Introducción de Circuit Breaker}: mediante Resilience4j
(\texttt{@CircuitBreaker(name = "catalogo", fallbackMethod = "...")}
sobre \texttt{CatalogoClient.obtenerProducto}), con tres estados:

\[
\text{CLOSED} \xrightarrow{\text{N fallos consecutivos}} \text{OPEN}
\xrightarrow{\text{cooldown}} \text{HALF-OPEN}
\begin{cases}
\text{éxito} \to \text{CLOSED} \\
\text{falla} \to \text{OPEN}
\end{cases}
\]

La diferencia clave frente al manejo actual (\texttt{try/catch} + timeout):
en estado \texttt{OPEN}, las peticiones fallan \emph{inmediatamente} sin
esperar los 2 segundos de timeout ni consumir un hilo del \emph{pool} HTTP por
intento —protegiendo a \texttt{svc-pedidos} de una falla en cascada cuando
\texttt{svc-catalogo} está caído y el sistema recibe tráfico sostenido.

% ======================================================
\section{Conclusiones}
% ======================================================

\begin{enumerate}[leftmargin=2em]
    \item El prototipo demuestra los cinco pilares exigidos por la Unidad 2:
    separación de dominios con \emph{Database per Service}, diseño RESTful de
    Nivel 2 de Richardson, comunicación inter-servicio síncrona con manejo
    explícito de fallos, contenerización con \emph{multi-stage builds}, y
    orquestación reproducible con un solo comando (\texttt{docker compose up
    --build}).
    \item El acoplamiento temporal identificado en la comunicación síncrona
    REST es una limitación consciente del alcance del examen; en producción
    se mitigaría con un modelo orientado a eventos, tal como se discute en la
    Sección~\ref{sec:analisis}.
    \item La arquitectura implementada es exactamente la base sobre la que se
    construirá el Proyecto Fin de Curso, validando que las decisiones de
    diseño tomadas aquí (API Gateway, Database per Service, JWT pendiente de
    integrar) son extensibles.
\end{enumerate}

% ======================================================
\begin{thebibliography}{9}
% ======================================================

\bibitem{richardson2018}
C.~Richardson, \emph{Microservices Patterns: With Examples in Java}.
Shelter Island, NY, USA: Manning Publications, 2018.

\bibitem{fielding2000}
R.~T.~Fielding, ``Architectural Styles and the Design of Network-based
Software Architectures,'' Ph.D.\ dissertation, Univ.\ of California, Irvine,
Irvine, CA, USA, 2000.

\bibitem{newman2021}
S.~Newman, \emph{Building Microservices: Designing Fine-Grained Systems},
2nd~ed. Sebastopol, CA, USA: O'Reilly Media, 2021.

\bibitem{aksakalli2021}
I.~Karabey~Aksakalli, T.~Çelik, A.~B.~Can, and B.~Tekinerdoğan,
``Deployment and communication patterns in microservice architectures: A
systematic literature review,'' \emph{Journal of Systems and Software},
vol.~180, art.~111014, 2021, doi: 10.1016/j.jss.2021.111014.

\bibitem{blinowski2022}
G.~Blinowski, A.~Ojdowska, and A.~Przybyłek, ``Monolithic vs.\ microservice
architecture: A performance and scalability evaluation,'' \emph{IEEE Access},
vol.~10, pp.~20357--20374, 2022, doi: 10.1109/ACCESS.2022.3152803.

\end{thebibliography}

\end{document}
